\documentclass[11pt,a4paper]{article}
\usepackage[margin=2.2cm]{geometry}
\usepackage[T1]{fontenc}
\usepackage[utf8]{inputenc}
\usepackage{lmodern}
\usepackage{booktabs}
\usepackage{longtable}
\usepackage{array}
\usepackage{tabularx}
\usepackage{enumitem}
\usepackage{fancyhdr}
\usepackage{titlesec}
\usepackage{hyperref}
\usepackage{xcolor}
\usepackage{setspace}
\usepackage{microtype}

\hypersetup{
  colorlinks=true,
  linkcolor=blue!50!black,
  urlcolor=blue!60!black,
  citecolor=blue!50!black,
  pdftitle={Open Clinical Record --- Software Requirements Specification v2.0},
  pdfauthor={Open Clinical Record Internship Project}
}

\pagestyle{fancy}
\fancyhf{}
\fancyhead[L]{\small Open Clinical Record}
\fancyhead[R]{\small SRS v2.0}
\fancyfoot[C]{\thepage}
\renewcommand{\headrulewidth}{0.4pt}

\titleformat{\section}{\large\bfseries\color{blue!40!black}}{\thesection}{0.8em}{}
\titleformat{\subsection}{\normalsize\bfseries}{\thesubsection}{0.7em}{}
\setlist[itemize]{leftmargin=1.4em,itemsep=0.2em}
\setlist[enumerate]{leftmargin=1.6em,itemsep=0.2em}
\onehalfspacing

\begin{document}

\begin{titlepage}
\centering
\vspace*{1.5cm}
{\LARGE\bfseries Open Clinical Record}\\[0.4cm]
{\Large Software Requirements Specification}\\[1.2cm]
{\large Version 2.0 --- Week 8 Finalization Baseline}\\[0.6cm]
{\large 27 September 2026}\\[1.5cm]
\begin{tabular}{ll}
\toprule
\textbf{Product} & Open Clinical Record (OCR) outpatient EMR \\
\textbf{Document type} & Software Requirements Specification \\
\textbf{Status} & Canonical requirements baseline \\
\textbf{Repository} & \texttt{elan-thinks/open-clinical-record} \\
\textbf{Branch} & \texttt{main} \\
\bottomrule
\end{tabular}
\vfill
{\small Internship MVP --- Patient Management, Patient Chart, Appointment Management}\\[0.3cm]
{\small This document prevails over earlier drafts where they conflict.}
\end{titlepage}

\tableofcontents
\newpage

\section{Introduction}

\subsection{Purpose}
This Software Requirements Specification (SRS) defines the functional and non-functional requirements for the Open Clinical Record internship Minimum Viable Product (MVP). It is the authoritative requirements baseline for design validation, implementation, testing, demonstration, and formal evaluation.

Where earlier requirement drafts, research notes, or UI mocks conflict with this document, \textbf{this version prevails}.

\subsection{Scope}
Open Clinical Record is a focused \textbf{outpatient} electronic medical record supporting three business modules only:

\begin{enumerate}
  \item \textbf{Patient Management} --- registration, unique identification, search, demographic updates, and patient status (including deceased policy).
  \item \textbf{Patient Chart} --- longitudinal clinical context: allergies, medical history items, and \textbf{per-visit} documentation (vital signs, diagnoses, clinical notes), without overwriting prior attendances.
  \item \textbf{Appointment Management} --- scheduling, rescheduling, cancellation with reason, status transitions, check-in/queue, and appointment event history.
\end{enumerate}

The product is \textbf{not} a full hospital EMR. Explicitly excluded: pharmacy dispensing, laboratory and imaging order management, billing and insurance, patient portal and native mobile apps, FHIR/HL7 external exchange, multi-facility enterprise scheduling, and AI clinical decision support.

\subsection{Definitions}
\begin{tabularx}{\textwidth}{@{}>{\bfseries}p{3.2cm}X@{}}
\toprule
Term & Definition \\
\midrule
MVP & Minimum Viable Product for the internship deliverable \\
Patient & Registered person with a unique medical record number (MRN) \\
Appointment & Scheduled (or walk-in typed) attendance slot with status lifecycle \\
Visit / ClinicalVisit & One facility attendance record holding clinical content \\
Draft & Editable visit status \\
Final & Closed visit; clinical content is immutable \\
JWT & JSON Web Token used for API authentication \\
RBAC & Role-based access control enforced at the API \\
\bottomrule
\end{tabularx}

\subsection{References}
Project scope, clinical domain rules, non-functional requirements, traceability matrix, access-control report, clinical visit model, technical documentation, and user guide in the project repository under \texttt{docs/}.

\subsection{Document overview}
Section~2 describes the product and users. Section~3 states functional requirements. Section~4 summarises non-functional requirements. Section~5 covers interfaces and data. Section~6 defines verification and acceptance. Section~7 lists exclusions.

\section{Overall Description}

\subsection{Product perspective}
OCR is a standalone web application with a clear three-tier structure:

\begin{verbatim}
  Browser (React SPA + TypeScript)
           |  REST / JSON + JWT
  ASP.NET Core 8 Web API
           |  EF Core 8 + Npgsql
  PostgreSQL
\end{verbatim}

Controllers remain thin; application services encapsulate workflows (patients, appointments, clinical chart, dashboard, audit).

\subsection{Product functions}
\begin{itemize}
  \item Authenticate users and authorize operations by role at the API boundary.
  \item Register and manage patients, including deceased mark/clear policy.
  \item Maintain a longitudinal chart and per-visit clinical documentation.
  \item Schedule and manage appointments; check in patients; preserve history.
  \item Provide role-aware dashboards and basic operational statistics.
  \item Audit important actions for administrative review.
\end{itemize}

\subsection{User characteristics --- application roles}
The MVP defines \textbf{exactly four} application roles:

\begin{tabularx}{\textwidth}{@{}p{3.4cm}p{2.4cm}X@{}}
\toprule
\textbf{Role} & \textbf{Claim} & \textbf{Primary responsibilities} \\
\midrule
Receptionist / Front Desk & \texttt{Receptionist} & Registration, booking, check-in, cancel/reschedule, mark deceased \\
Nurse / Clinical Staff & \texttt{Nurse} & Chart support, vitals, visit documentation (write) \\
Clinician / Doctor & \texttt{Doctor} & Chart review, clinical documentation, mark/clear deceased \\
System Administrator & \texttt{Admin} & Users, audit, configuration; operational access without default clinical authorship \\
\bottomrule
\end{tabularx}

\vspace{0.4em}
\noindent\textbf{Security principle:} UI visibility is not security. The API shall enforce authorization independently of the client.

\subsection{Constraints}
\begin{itemize}
  \item Internship time-box limits scope to the three modules above.
  \item PostgreSQL is the system of record; EF Core migrations are schema authority.
  \item Development training credentials must never be used in production.
  \item International interoperability standards (e.g.\ FHIR) are not MVP acceptance criteria.
\end{itemize}

\subsection{Assumptions and dependencies}
Users have a modern browser and network access to the API. PostgreSQL is available for non-test deployments. Appointment dates are treated as clinic-local calendar dates.

\section{Functional Requirements}

Priority \textbf{M} means mandatory for MVP acceptance. Identifiers align with the requirements traceability matrix.

\subsection{Patient management (FR-PM)}

\begin{longtable}{@{}p{2.2cm}p{11.2cm}c@{}}
\toprule
\textbf{ID} & \textbf{Requirement} & \textbf{Pri} \\
\midrule
\endhead
FR-PM-001 & Authorized users shall register a new patient with required demographics. & M \\
FR-PM-002 & The system shall assign a unique medical record number to each patient. & M \\
FR-PM-003 & Authenticated staff shall search patients by name and/or MRN. & M \\
FR-PM-004 & Authorized users shall view a patient profile. & M \\
FR-PM-005 & Authorized users shall update permitted demographic and contact fields. & M \\
FR-PM-006 & The system should support duplicate awareness (e.g.\ search-before-create). & M \\
FR-PM-007 & The system shall maintain patient status: Active, Inactive, or Deceased. & M \\
FR-PM-008 & Marking a patient deceased (with provenance) shall be allowed for Admin, Doctor, and Receptionist --- not Nurse. & M \\
FR-PM-009 & Clearing deceased status shall be allowed for Admin and Doctor only. & M \\
FR-PM-010 & The system shall reject new appointments for patients with status Deceased. & M \\
\bottomrule
\end{longtable}

\subsection{Patient chart and visits (FR-PC)}

\begin{longtable}{@{}p{2.2cm}p{11.2cm}c@{}}
\toprule
\textbf{ID} & \textbf{Requirement} & \textbf{Pri} \\
\midrule
\endhead
FR-PC-001 & Authorized users shall open a patient chart in single-patient context. & M \\
FR-PC-002 & The system shall display a chart summary (identity, status, key alerts). & M \\
FR-PC-003 & Doctor and Nurse shall record and view allergies. & M \\
FR-PC-004 & Doctor and Nurse shall record and view medical history items. & M \\
FR-PC-005 & Clinical attendance shall be represented as a ClinicalVisit linked to the patient. & M \\
FR-PC-006 & Visit statuses shall include Draft, Final, and Cancelled. & M \\
FR-PC-007 & New visits shall default to Draft and remain editable until Final. & M \\
FR-PC-008 & Final visits shall be immutable (no further clinical documentation edits). & M \\
FR-PC-009 & Clinical staff shall record vital signs on a visit. & M \\
FR-PC-010 & Clinical staff shall record diagnoses (primary/secondary) on a visit. & M \\
FR-PC-011 & Clinical staff shall record clinical notes and plan/instructions on a visit. & M \\
FR-PC-012 & Prior visits shall never be overwritten; a return creates a new visit. & M \\
FR-PC-013 & Visit history shall be displayed chronologically on the chart. & M \\
FR-PC-014 & Receptionist shall not create or edit clinical visit content (API forbidden). & M \\
FR-PC-015 & Finalizing a visit linked to an appointment shall set that appointment to Completed. & M \\
\bottomrule
\end{longtable}

\subsection{Appointments (FR-AP)}

\begin{longtable}{@{}p{2.2cm}p{11.2cm}c@{}}
\toprule
\textbf{ID} & \textbf{Requirement} & \textbf{Pri} \\
\midrule
\endhead
FR-AP-001 & Authorized staff shall create an appointment for an existing non-deceased patient. & M \\
FR-AP-002 & Appointments shall be viewable for a selected date (list). & M \\
FR-AP-003 & A calendar view (day/week) of appointments shall be supported. & M \\
FR-AP-004 & Appointment types shall be supported as labels (Consultation, Follow-up, Walk-in, etc.). & M \\
FR-AP-005 & Reschedule shall be allowed for Scheduled, Waiting, Cancelled, NoShow; blocked for CheckedIn, InProgress, Completed. & M \\
FR-AP-006 & Cancel shall require a non-empty reason. & M \\
FR-AP-007 & Statuses Scheduled, Waiting, CheckedIn, InProgress, Completed, Cancelled, NoShow shall have enforced transitions. & M \\
FR-AP-008 & Check-in shall transition to CheckedIn and create a Draft visit when applicable. & M \\
FR-AP-009 & Meaningful changes shall record AppointmentEvent history. & M \\
FR-AP-010 & The active check-in queue shall exclude Completed, Cancelled, and NoShow. & M \\
FR-AP-011 & Booking shall be available to Admin, Doctor, Nurse, and Receptionist. & M \\
\bottomrule
\end{longtable}

\subsection{Security, validation, and audit}

\begin{longtable}{@{}p{2.2cm}p{11.2cm}c@{}}
\toprule
\textbf{ID} & \textbf{Requirement} & \textbf{Pri} \\
\midrule
\endhead
FR-SEC-001 & Users shall authenticate before accessing protected resources. & M \\
FR-SEC-002 & Exactly four application roles shall be enforced as defined above. & M \\
FR-SEC-003 & Unauthorized operations shall be rejected at the API (401/403). & M \\
FR-SEC-004 & Important actions shall be auditable with actor and timestamp. & M \\
FR-SEC-005 & Admin shall be able to list audit events. & M \\
FR-VAL-001 & Required fields and relationships shall be validated before persistence. & M \\
FR-VAL-002 & Patient and appointment business status rules shall be enforced. & M \\
FR-VAL-003 & Error messages shall be clear without leaking sensitive internals. & M \\
FR-RPT-001 & Dashboard statistics shall be available (e.g.\ today, waiting, checked-in). & M \\
FR-RPT-002 & A simple reports view based on live stats may be provided (not a BI suite). & M \\
\bottomrule
\end{longtable}

\section{Non-Functional Requirements (Summary)}

Detailed NFRs are maintained in \texttt{docs/03-requirements/non-functional-requirements.md}. MVP commitments include:

\begin{itemize}
  \item \textbf{Security:} backend authorization, hashed passwords, least privilege, no secrets in source control.
  \item \textbf{Privacy:} minimum necessary access; synthetic data in demonstrations.
  \item \textbf{Integrity:} unique MRN; consistent relationships; history not silently deleted.
  \item \textbf{Performance:} common reads typically within approximately 2\,s and writes within approximately 3\,s under normal laboratory load; target of at least ten concurrent authenticated users.
  \item \textbf{Reliability:} failed operations must not report success; confirmed data preserved.
  \item \textbf{Usability:} role-aware navigation; visible patient context; distinguishable appointment statuses.
  \item \textbf{Auditability:} actor and timestamp; audit stream not editable through ordinary workflows.
  \item \textbf{Maintainability:} service-layer separation; migrations as schema source of truth.
\end{itemize}

\section{Interfaces and Data}

\subsection{User interfaces}
Web SPA with role-aware shell. Core screens: Login, Dashboard, Patients, Register, Chart, Appointments (list and calendar), Check-in queue, Profile, Admin audit.

\subsection{Software interfaces}
REST JSON API under \texttt{/api/*}; PostgreSQL via EF Core~8.

\subsection{Logical data}
Patient (status, unique MRN); allergies and history items; ClinicalVisit with vital signs, diagnoses, notes; Appointment and AppointmentEvent; PatientDeathRecord; AuditEvent; identity tables for users and roles. \textbf{EF Core migrations} are authoritative; reference SQL under \texttt{docs/05-data/} is secondary.

\section{Verification and Acceptance}

\subsection{Verification methods}
\begin{itemize}
  \item Automated API tests: access matrix, appointment transitions, clinical documentation/finalize, longitudinal visits.
  \item Manual end-to-end checklist under \texttt{docs/09-testing/}.
  \item Continuous integration: backend tests and frontend typecheck/build on \texttt{main}.
\end{itemize}

\subsection{MVP acceptance criteria}
\begin{enumerate}
  \item Four roles can sign in with role-appropriate navigation.
  \item A patient can be registered, searched, and opened in the chart.
  \item An appointment can be booked, checked in, cancelled with reason, and rescheduled under allowed statuses.
  \item Clinical staff can document a Draft visit and finalize; history shows multiple visits without overwrite.
  \item Finalizing a visit completes the linked appointment for queue purposes.
  \item Unauthorized clinical writes by Receptionist are rejected.
  \item Deceased booking is rejected; mark/clear deceased follows the role matrix.
  \item Technical documentation and user guide are available in the repository.
\end{enumerate}

\section{Exclusions}
Prescription and pharmacy, laboratory and imaging workflows, billing, patient portal, FHIR exchange, multi-facility enterprise scheduling, and AI decision support are \textbf{not} MVP requirements.

\section*{Document control}
Version 2.0 (27 September 2026). Implementation evidence is the \texttt{main} branch of the project repository. Subsequent changes shall increment the version and date.

\end{document}
